\subsection{Arbitrary products of the noncomplete henselian families}
\label{C:henselian-products-subsection}

In this subsection a \emph{solved local factor} means one of
Theorem~\ref{C:effective-henselian-DVR-WG},
Theorem~\ref{C:finite-field-henselization-WG}, or
Theorem~\ref{C:arithmetic-henselization-WG}, with its supplied
finite input presentation.  Thus all its canonical indices,
actual generators and operations are given by ordinary
terminating algorithms.  The proofs apply as well to any
additional henselian local family for which those same
finite-name and kernel algorithms have been supplied.
We retain the notation of (C.3.HD2).

\begin{proposition}[Proved: exact minimal ranks and coefficient bounds]
\label{C:henselian-minimal-form-bounds}
Let $R$ be a solved local factor.  Write an element $\xi$ in
its canonical form coordinates as
\[
 \xi=\sum_{j\in\mathscr I_{\rm s}}(a_j e_j+b_j t_j)
       +\sum_{j\notin\mathscr I_{\rm s}}a_j e_j,
 \qquad a_j\in\mathbf Z,\quad b_j\in\{0,1\},             \tag{C.3.HP1}
\]
with finite support.  Let $n_j$ and $d_j$ be the free rank
and minimal coefficient degree of the indexed indecomposable
$M_j$, and put
\[
 w_j=\begin{cases}
 n_j,&j\in\mathscr I_{\rm s}\cup\mathscr I_{\rm u},\\
 2n_j,&j\in\mathscr I_{\rm a}\cup\mathscr I_{\rm h},
 \end{cases}
 \qquad
 \delta_j=\begin{cases}
 1,&j\in\mathscr I_{\rm s},\ a_j=0,\ b_j=1,\\
 0,&\text{otherwise}.
 \end{cases}
\]
Define
\[
 \begin{split}
 r_+(\xi)&=\sum_jw_j\max(a_j,0)
                       +\sum_{j\in\mathscr I_{\rm s}}n_j\delta_j,\\
 r_-(\xi)&=\sum_jw_j\max(-a_j,0)
                       +\sum_{j\in\mathscr I_{\rm s}}n_j\delta_j,\\
 c(\xi)&=\max\bigl(\{d_j:a_j\ne0\text{ or }
                      (j\in\mathscr I_{\rm s},b_j=1)\}\cup\{0\}\bigr).
 \end{split}                                             \tag{C.3.HP2}
\]
For every expression $\xi=[P,q]-[Q,h]$ by two actual
symmetric representations,
\[
 \operatorname{rank}P\geq r_+(\xi),\qquad
 \operatorname{rank}Q\geq r_-(\xi),\qquad
 \max(\deg P,\deg Q)\geq c(\xi).
\]
All three bounds are attained simultaneously by a computable
pair of actual forms.  In particular $\xi$ is such a difference
with both ranks at most $B$ and both coefficient degrees at
most $D$ if and only if
\[
                    r_+(\xi),r_-(\xi)\leq B,
                    \qquad c(\xi)\leq D.                \tag{C.3.HP3}
\]
This is an ordinary decidable test from its finite coordinates.
\end{proposition}
\begin{proof}
In a symmetric block the monoid of actual forms consists
of pairs $(r,b)$ with $r\geq0$, $b\in\mathbf Z/2$, subject
to $b=0$ when $r=0$.  This follows from the proved finite
residue-field form classification and henselian lifting.
The underlying module is $M_j^r$, and its class is
$r e_j+b t_j$.  If the positive and negative forms have
block multiplicities $u,v$, then $u-v=a_j$ and their
bits differ by $b_j$.
For $a_j>0$ the smallest possible multiplicities are
$(u,v)=(a_j,0)$; the positive form may carry either
required bit.  For $a_j<0$ they are $(0,-a_j)$, with
that bit on the negative form; subtraction does not
change a bit of order two.  For $a_j=0,b_j=1$, both
multiplicities must be positive, so their simultaneous
minimum is $(1,1)$.  For $a_j=b_j=0$ it is $(0,0)$.
These are exactly the symmetric contributions in
(C.3.HP2), and are individually minimal on both sides.

In each of the other types, actual forms have one
nonnegative integer coordinate.  The generator rank
is $w_j$: it is $n_j$ for unitary type and $2n_j$
for an alternating or dual-pair hyperbolic generator.
Their smallest positive and negative multiplicities
are $\max(a_j,0)$ and $\max(-a_j,0)$.
Krull--Schmidt multiplicities and the residue-form
classification are independent between the blocks.
Summing the individual lower bounds therefore proves
the two rank inequalities, and taking the indicated
orthogonal sums attains them both.

A representation whose coefficients lie in $C_D(R)$
has all its direct summands in $C_D(R)$: choose free
bases of the summands, and conjugate the coaction by
one constant invertible matrix over $R$.
Every block coefficient remains an $R$-linear
combination of the original coefficients.
Thus the minimal degree of every occurring
indecomposable is at most $D$.
Each nonzero coordinate in (C.3.HP1) forces its
indecomposable, or its dual pair, to occur in at
least one of $P,Q$.  Duals have the same minimal
coefficient degree, since the antipode preserves
$C_d$ and its square is the identity.
This proves the lower bound $c(\xi)$.

The explicit forms attaining the rank minima use
only those supported indecomposables.  In a nonzero
symmetric block, a requested bit is represented by
one copy of the nonsquare form $q_j a_j'$ and the
remaining copies of $q_j$; for the isolated torsion
coordinate use $(q_j a_j')-(q_j)$.
Here $a_j'$ denotes the previously constructed
self-adjoint nonsquare lift, not the integer
coefficient $a_j$.
Changing the invariant form never changes its
coaction.  The unitary and hyperbolic generators
also have coefficient degree $d_j$.
Hence the same pair has common degree at most
$c(\xi)$, and all three minima are attained.
The zero element uses two zero modules, with degree
convention zero.  All input ranks, degrees, integers
and bits are computable, which proves (C.3.HP3).
\end{proof}

\begin{theorem}[Proved: exact canonical coordinates for arbitrary henselian products]
\label{C:henselian-products-WG}
Let $(R_i)_{i\in I}$ be any family of solved local factors,
and put $R=\prod_{i\in I}R_i$.
There is an injective homomorphism preserving products
and every exterior operation,
\[
 WG_{\rm sp}(BO_{m,R})\longrightarrow
                 \prod_{i\in I}WG_{\rm sp}(BO_{m,R_i}).   \tag{C.3.HP4}
\]
In the literal canonical coordinates of the local factors,
its image is exactly the set of families $(\xi_i)$ for
which the following three nonnegative integer functions
are bounded on $I$:
\[
              i\longmapsto r_+(\xi_i),\qquad
              i\longmapsto r_-(\xi_i),\qquad
              i\longmapsto c(\xi_i).                    \tag{C.3.HP5}
\]
Each value in (C.3.HP5) is obtained by a finite calculation
on that coordinate.  Products and exterior operations in
this presentation are coordinatewise and are computed
by the ordinary local algorithms at every requested
coordinate.  No finite algorithm deciding equality or
boundedness of an arbitrary infinite family is asserted.

This includes infinite non-Noetherian products of the
noncomplete rings $(\mathbf F_q[t]_{(t)})^h$.
For example the countable product of one such ring
admits no profinite ring topology, so this case is not
contained in the preceding profinite classification.
\end{theorem}
\begin{proof}
Every finite projective $R$-module is the image of
an idempotent in a finite matrix algebra over $R$.
Its rank at every local factor is bounded by that
ambient matrix size.  Its finite coaction coefficients
lie in one $C_D(R)$, and reduction to each factor
preserves that bound.  Thus any difference of two
actual forms maps to a family satisfying (C.3.HP5),
by Proposition~\ref{C:henselian-minimal-form-bounds}.

Suppose a difference $[P,q]-[Q,h]$ maps to zero in
every factor.  Henselian residue-form classification
has cancellation: each local monoid is the product
of the cancellative block monoids used in the
preceding proposition.  Hence there is an actual
equivariant isometry $P_i\simeq Q_i$ for each $i$.
Write $P=eR^a$ and $Q=fR^b$.
The local isometry and its inverse can be represented
by matrices
\[
 g_i=f_i g_i e_i\in M_{b\times a}(R_i),
 \qquad u_i=e_i u_i f_i\in M_{a\times b}(R_i).
\]
Their fixed matrix sizes permit their entries to
be assembled in the product ring, yielding
$g=fge$ and $u=euf$ with $ug=e$ and $gu=f$.
The Gram and equivariance equations hold coordinatewise.
For equivariance, the finite free integral coefficient
basis makes the map from any fixed $C_D(R)$ to
$\prod_i C_D(R_i)$ injective, so these equations
also hold in the actual Hopf algebra over $R$.
Thus $g$ is an actual equivariant isometry, proving
injectivity in (C.3.HP4).  No global cancellation
hypothesis on the product ring was used.

Conversely assume (C.3.HP5), and choose common
bounds $B,D$.  The preceding proposition constructs
actual local forms $(P_i,q_i),(Q_i,h_i)$ with
common rank bound $B$ and degree bound $D$, whose
difference is $\xi_i$.
Choose their specified free bases.  For example,
if $r_i=\operatorname{rank}P_i$, let
$e_i=\operatorname{diag}(I_{r_i},0)$ in $M_B(R_i)$.
Pad its coaction and Gram matrix with zeros on the
complement, as in
Lemma~\ref{C:profinite-projective-data}; the padded
coaction has counit $e_i$, not $I_B$.
Pad the inverse Gram matrix with zeros as well.

The idempotents assemble to $e\in M_B(R)$.
The coaction arrays assemble because $C_D$ is a
fixed finite free integral coefficient module:
there are only finitely many coefficient positions,
and the values at each such position define one
element of $R$.  The padded Gram matrices and their
inverse maps assemble in the same manner.
All idempotent, coaction, invariance and inverse-Gram
identities are finite coefficient identities and
hold in every coordinate, hence hold over $R$.
They define a genuine symmetric rational representation
on $eR^B$ with a nonsingular form.  Perform the
same construction for the negative forms.
Their difference maps to $(\xi_i)$, proving the
exact image assertion.  Variable local ranks were
absorbed by the idempotents; a global free-module
assumption was not made.

Base change preserves tensor products and exterior
forms, so operations in (C.3.HP4) are coordinatewise.
They preserve the image condition.  Concretely, for
two differences with common rank bounds $B,B'$
and degree bounds $D,D'$, multiplication is a
difference of two actual forms with rank bound
$2BB'$ and degree bound $D+D'$.
For a specified exterior coefficient $r$, expand
$\lambda_t(P)\lambda_t(Q)^{-1}$ to order $r$.
It is a finite integral combination of tensor
products of exterior forms, with a rank bound
in terms of $B,r$ and degree at most $rD$.
Grouping positive and negative terms supplies
common bounds again.  Their minimal bounds can
then be read from (C.3.HP2) at each coordinate.
The local algorithms compute every such coordinate,
which is exactly the access claimed in the theorem.

Finally put $R_0=(\mathbf F_q[t]_{(t)})^h$ and
$R=R_0^{\mathbf N}$.
The ideals generated by the first $n$ coordinate
idempotents form a strictly increasing chain,
so $R$ is not Noetherian.
If $R$ had a profinite ring topology, the image
of multiplication by its first coordinate
idempotent would be a closed compact subring
isomorphic as a ring to $R_0$.
But $R_0$ is countably infinite, whereas an infinite
profinite additive group is uncountable.
For the latter fact, an infinite compact Hausdorff
group has no isolated points; recursively choose
two disjoint nonempty open subsets with disjoint
closures inside each previously chosen open set.
Compactness gives a different point for each
infinite binary branch, proving uncountability.
This contradiction proves the last assertion.
\end{proof}
